% Generated by task tex-groundtruth from dossiers/gt-bus.md (section G, findings G1 to G17),
% with the corrections of dossiers/verify-gt-bus.md.
\section{The setup, the commitment and the bus phase: findings}\label{sec:gen-bus-findings}

\begin{sloppypar}
Short forms in this section. Specification files are under \file{doc/leanvm/body/}:
\code{05:} is \file{05-arithmetization.tex}, \code{08:} is \file{08-end-to-end-protocol.tex}.
Rust files are under \file{crates/lean_vm/src/}: \code{mod.rs} is \file{cpu/mod.rs},
\code{layout.rs} is \file{cpu/layout.rs}, and \code{leaf.rs}, \code{gkr.rs},
\code{constraints.rs} are at that root; \code{fs/} is \file{crates/fiat_shamir/src/}. \code{py:} is
\file{python-verifier/verifier.py}. These are at leanVM's pin \code{a386121f}. \code{bp:} is the
blueprint \file{docs/roadmap/protocol-blueprint.md} and \code{st:} the status
\file{docs/roadmap/protocol-status.md}, both at leanerVM \code{b435631} (the merge of the Lean
4.34.1 upgrade moved blueprint lines after its line 51 by four). \code{Phase.lean},
\code{Seams.lean}, \code{Instance.lean} are under \file{LeanerVM/Protocol/Spine/};
\code{Component.lean} and \code{KnowledgeAppend.lean} under \file{LeanerVM/Protocol/ToArkLib/};
\code{Statement.lean} is \file{LeanerVM/Arithmetization/Statement.lean}; all at \code{b435631}.
ArkLib is cited at its old pin \code{dca90385}. The number after the colon is the line.
\end{sloppypar}

% ---------------------------------------------------------------- G1
\begin{finding}{major}{f:bus-gkr-normalized}{The GKR layer sumcheck is the normalized one: the blueprint gives it the wrong degree, the wrong error and no definition}
\fhead{Classification} An error of the blueprint.

\fhead{Evidence} The deployed round message is the cofactor of degree 4, four elements on the
wire, $c_0$ derived through the equality factor; the running claim is the cofactor's value; the
layer's final check has no equality factor (\cref{sec:gen-bus-gkr}; \code{gkr.rs:399-404, 410-413};
\code{fs/transcript.rs:297-302}; \code{py:411-413, 445, 449}). The probe's honest prover, written to
that reading, is accepted by the pinned Python verifier (dossier \file{gt-bus.md}, section H.1;
the citation check re-ran the probe and obtained the same output). The table sumcheck is the
other variant: three of the cubic's four coefficients are sent, $c_1$ derived, and the verifier
multiplies the equality factors into its final check (\code{constraints.rs:267}
\code{next_round_poly(4, claim, None)}, \code{:271-274}; \code{py:608-614}). (The citation check
corrected the dossier's ``the whole cubic is sent, $c_1$ derived'', which contradicts itself:
the wire carries $c_0, c_2, c_3$.)

The blueprint:
\begin{itemize}
\item \code{bp:321}: ``Layer sumcheck messages have degree 5 (eq × four multilinears); the round
  check is on the cofactor (Gruen).''
\item \code{bp:315}: ``A round polynomial travels as its coefficient list, low degree first, of
  length $d + 1$; the oracle protocol sends all of them. Dropping one is the \emph{encoding} of
  Layer 12.''
\item \code{bp:898-901}: the eq-weighted variant has ``$\eq(ζ_{<τ_j}, ·)$ as an explicit factor
  whose evaluation at $r$ the verifier computes itself (so $d$ counts the cofactor plus one)''.
\item \code{bp:941-942}: ``$5/|\E|$ to each sumcheck round''. The tracker, in its hole for the
  generic GKR (G5): ``degree-5 round polynomials''.
\end{itemize}
Read literally, a message of degree 5 has six coefficients and five on the wire, where the Rust
has four: a verifier built so rejects every Rust proof. Read with Layer 4's variant, the verifier
multiplies the equality factors in, which is a different verifier from the deployed one (it
accepts when $χ_t = 1 + \mathit{point}[t]$, whatever the children). Layer 12's
\code{RoundPoly.decode (d) (claim) (eq? : Option E)} (\code{bp:1212}) knows both encodings;
Layers 4 and 5 define only one protocol.

\fhead{Proposed change}
\begin{change}{the GKR row of the conventions (\code{bp:321})}
\asis ``Layer sumcheck messages have degree 5 (eq × four multilinears); the round check is on
the cofactor (Gruen).''
\asproposed ``The layer sumcheck is normalized: the prover's message in round $t$ is the
cofactor $h_t$, of degree 4 (the four children's lines multiplied, summed against the equality
kernel of the coordinates not yet bound), five coefficients; the verifier checks
$(1 + r_t)·h(0) + r_t·h(1)$ against the running claim, $r_t$ the coordinate of the layer's
point; the new running claim is $h(χ_t)$; the layer's final check compares the running claim
with $Σ_s λ^s ∏ \mathit{children}$, with no equality factor (\code{gkr.rs:399-413}). The table
sumcheck is not normalized (\code{constraints.rs:267-274}).''
\reason Faithfulness of the wire format and of the verdict; the error is then tight.
\end{change}
Also:
\begin{itemize}
\item Layer 4: define the two variants by name, the plain one (error $d/|F|$, $d$ the degree of
  the polynomial sent, the verifier evaluating the equality factors at the end) and the
  normalized one (error $d_c/|F|$, $d_c$ the degree of the cofactor).
\item \code{bp:941-942} and the tracker: ``$4/|\E|$ to each sumcheck round''.
\end{itemize}
\end{finding}

% ---------------------------------------------------------------- G2
\begin{finding}{major}{f:bus-extra-combiner}{One combiner more than layers: the last is drawn and never used}
\fhead{Severity} Major, by the brief's definition: the blueprint omits a challenge the verifiers
draw; the fix is one sentence.

\fhead{Classification} An error of the blueprint.

\fhead{Evidence} \code{gkr.rs:369} (before the loop) and \code{gkr.rs:391, 423} (at the end of
every iteration, the last included); \code{py:435, 453}. The dossier's probe (section H.1): the
GKR draws $μ²/4 + μ + 1$ challenges for even $μ$, and $ζ$ never contains the last. The
specification has ``at every layer a combiner'' (\code{05:91}). The blueprint: ``A fresh combiner
$λ$ per layer'' (\code{bp:321}), and Layer 5 ``a fresh $λ$ per layer … exactly as pinned''
(\code{bp:940-941}).

In the oracle model a challenge nobody uses is harmless. After Fiat--Shamir it is a squeeze that
changes the chain's state (\code{fs/lib.rs:95-99}), so $ξ$ and everything after depend on it.

\fhead{Proposed change}
\begin{change}{the GKR row of the conventions (\code{bp:321})}
\asis ``A fresh combiner $λ$ per layer''.
\asproposed ``A combiner $λ$ is drawn after the roots and after every layer, the last layer
included; that last one is used by nothing and its error is $0$ (\code{gkr.rs:369, 423}).''
\end{change}
Add the message count of the bus phase (the element counts after \cref{tab:gen-bus-phase}) to
acceptance test 17.
\end{finding}

% ---------------------------------------------------------------- G3
\begin{finding}{major}{f:bus-generic-gkr}{The generic GKR of Layer 5 cannot be appended in the bus phase, and cannot carry the zerocheck clause}
\fhead{Classification} An error of the blueprint; the first half is also a deviation an upstream
library forces.

\fhead{Evidence} Layer 5's component (\code{bp:930-937}):
\begin{leancode}
def gkr (nside μ : ℕ) : OracleReduction []ₒ
    (StmtIn := Fin nside → E)                                    -- the claimed roots
    (OStmtIn := fun _ : Fin nside ↦ ETable μ) Unit               -- the leaf tables
    (StmtOut := (Fin μ → E) × (Fin nside → E))                   -- ζ and the leaf claims Ṽ₀ˢ(ζ)
    (OStmtOut := fun _ : Fin nside ↦ ETable μ) Unit …
\end{leancode}
\src{docs/roadmap/protocol-blueprint.md:930-937 at b435631}

Its oracles are the leaf tables. A phase's one oracle is the stack (\code{Phase.lean:37-38}), and
two components are appended only when the output oracles of the first are the input oracles of
the second (\code{Component.lean:143-149}, \lean{Def.append}). The leaf tables are functions of
the stack and of $(α, β)$; carrying a reduction along such a map is ArkLib's context lifting,
whose completeness and knowledge-soundness theorems are admitted at the pin:
\begin{leancode}
-- .lake/packages/Arklib/ArkLib/OracleReduction/LiftContext/Reduction.lean:355-365, 542-560
theorem liftContext_completeness …  := by … sorry
theorem liftContext_rbr_knowledgeSoundness … := by … sorry
\end{leancode}
\src{ArkLib at dca90385, ArkLib/OracleReduction/LiftContext/Reduction.lean:355-365, 542-560 (the dossier's elided excerpt)}

The blueprint's ledger row on context lifting (A4) says of them ``Not consumed: statements are
shaped so that no lens is needed'' (\code{bp:262}). The holes table has the bus phase consume the
generic GKR holes (G5, G6) (\code{bp:606}), and the tracker writes
\code{busPhase I (G : Gkr.Def)}.

Second half. The bus phase's state function must turn \lean{ConstraintsVanish} into the claims
at $ζ$ while the coordinates of $ζ$ are drawn, which is inside the GKR's last layer (dossier
\file{gt-bus.md}, section D.3). The appended state function
(\lean{Verifier.KnowledgeStateFunction.appendGuarded}, \code{KnowledgeAppend.lean:474}, leanerVM's
port of ArkLib pull request \#615) is made of the components' own; a clause that is constant can
be framed around a component, a clause that changes with the component's challenges cannot. So
the security of a GKR that knows nothing of the constraints does not give the bus phase's. (The
citation check corrected the attribution: the dossier called it ArkLib's appended state
function, but at the pin \code{dca90385} ArkLib has no such declaration; the substance holds, the
definition is the components' own state functions spliced at the seam,
\code{KnowledgeAppend.lean:21-26, 474-488}.)

\fhead{Proposed change} Restate Layer 5's component over a context, with riders:
\begin{leancode}
-- proposed shape (Layer 5); S, OStmt arbitrary
def gkr (nside μ : ℕ) (leaves : S → (∀ i, OStmt i) → Fin nside → ETable μ)
    (riders : S → (∀ i, OStmt i) → List (Σ τ, ETable τ)) :
    Component.Def (S × (Fin nside → E)) OStmt Unit
                  (S × Vector E μ × (Fin nside → E)) OStmt Unit
-- relIn : the roots are the products of `leaves s o`, and every rider table is zero
-- relOut: the leaf claims hold at ζ, and every rider's extension vanishes at ζ_{<τ}
-- err   : 2/|E| per combiner (0 for the last), 4/|E| per round, 1/|E| per combination challenge
\end{leancode}
and say in Layer 6 that the bus phase is
\code{[(α, β); roots] ⟫ gkr ⟫ [boundary evaluations]} with \code{leaves} the three leaf tables of
the stack and \code{riders} the constraint tables. Replace \code{bp:262}'s ``no lens is needed'' by
the reason this shape needs none. The workaround is retired when ArkLib proves the lifting
theorems \textbf{and} offers a way to run a rider on a component's challenges; the first alone
does not give the second half.
\end{finding}

% ---------------------------------------------------------------- G4
\begin{finding}{major}{f:bus-admissible}{The admissibility predicate: the stacking window and the rate window}
\fhead{Classification} An error of the blueprint (three of its statements cannot hold together).

\fhead{Evidence} The verifiers reject $μ_{\mathrm{stack}} ∉ [15, 28]$ and $ρ ∉ [1, 4]$ (the rate
window and the stacking window, B8 and B9 of \cref{tab:gen-bus-checks}). The blueprint:
\begin{itemize}
\item \code{bp:800-801}: ``\code{def Sizes.Admissible (prog : Program) (s : Sizes) : Prop -- the caps; decidable}''
  and ``\code{theorem admissible_iff_caps : s.Admissible prog ↔ (Caps w ∧ Sizes.ofWitness w = some s)}''.
\item leanISA's \lean{Caps} has neither window: ``(7) the WHIR rate (\code{:167}) and (8) the
  stacked size \code{mu ∈ [MIN_MU, MAX_MU]} (\code{:174-176}) are parameters of the commitment,
  left to the protocol layer'' (\code{Statement.lean:79-81}; the structure is at \code{:250-263}).
\item \code{bp:1218-1220}: \lean{verify_iff_compiled} quantifies
  ``\code{∃ s, s.Admissible prog ∧ …}''.
\item \code{bp:1130}:
  \code{piopError_le (hs : s.Admissible prog) : Σ i, piopError s i ≤ 2^40/|E| + flockError}.
\end{itemize}
If \lean{Admissible} is the caps alone, \lean{verify_iff_compiled} is false (the Rust rejects
sizes that are admissible) and \lean{piopError_le} is false at the caps (dossier
\file{gt-bus.md}, section D.4). If it contains the windows, \lean{admissible_iff_caps} is false.
Besides, \code{Sizes.ofWitness w} cannot produce \code{logInvRate}, which no witness determines,
and \code{Caps w} has the conjunct \code{well_shaped}, which is no property of sizes.

\fhead{Proposed change} \code{bp:800-801}, as proposed:
\begin{leancode}
def Sizes.Admissible (prog : Program) (s : Sizes) : Prop   -- decidable; read_public, mod.rs:157-176
  -- κ_mem ∈ [16, 32] ∧ ∀ j, τ j ≤ 32 ∧ 3 ≤ τ 5 ∧ logInvRate ∈ [1, 4] ∧ μ_stack prog s ≤ 28
theorem caps_of_admissible (hs : s.Admissible prog) (hw : Sizes.heightsOf w = s.heights)
    (hd : WellShapedData w.data) : Caps w
\end{leancode}
with \lean{Sizes.ofWitness} returning the heights only. Layer 3's tests: add
$μ_{\mathrm{stack}} = 29$ and $\mathit{logInvRate} = 5$ rejected. State in
\lean{piopError_le}'s comment that the bound uses $μ_{\mathrm{bus}} ≤ 28$, a lemma of Layer 3
($N_{\mathrm{push}} ≤ Σ_{\mathrm{stack}} 2^κ$).
\end{finding}

% ---------------------------------------------------------------- G5
\begin{finding}{major}{f:bus-status-python-caps}{The status file's finding on the Python verifier's caps is false}
\fhead{Severity} Major, of the status file.

\fhead{Classification} An error of the status.

\fhead{Evidence} The status file's finding on the Python caps (F9), \code{st:316-318}: ``the
Python verifier omits the caps $\mathit{log\_mem} ∈ [16, 32]$, $τ_j ≤ 32$, the bytecode
power-of-two bound and $τ_{\mathrm{BLAKE2S}} ≥ 3$ (\code{verifier.py:1372-1379} versus
\code{cpu/mod.rs:158-170}): a divergence between the two verifiers, soundness-relevant, to report
upstream.'' \code{st:164}: ``Finding F9 … is still to be reported to leanVM.''

At the pin, \code{py:1378} calls \code{build_layout}, which begins:
\begin{srccode}
# python-verifier/verifier.py:856-864 at a386121f
def build_layout(bytecode: Sequence[K], log_memory: int, table_log_heights: Sequence[int]) -> Layout:
    log_bytecode = log2_strict(len(bytecode)) - BUS_BITS
    require(
        16 <= log_memory <= 32
        and all(0 <= log_height <= 32 for log_height in table_log_heights)
        and table_log_heights[OP_BLAKE2S] >= 3
        and 0 <= log_bytecode <= 32,
        "invalid announced table sizes",
    )
\end{srccode}
\src{python-verifier/verifier.py:856-864 at a386121f}

\code{log2_strict} requires a power of two (\code{py:252-254}). The check is present at the pin;
it was assembled in the commits \code{8a6e7750}, \code{348ca455} and \code{5751a5c7}, all
ancestors of \code{a386121f}. (The citation check corrected the dossier, which said the check was
introduced by commit \code{14fbca8f}: that commit's parent already has the same \code{require},
and \code{14fbca8f} only renamed the constant \code{BLAKE2S.opcode} to \code{OP_BLAKE2S}.) The Rust
itself says so: ``The other two verifiers reject it here too (\code{python-verifier},
\code{guests/aggregate.py})'' (\code{mod.rs:164-165}; the citation check corrected the line from
\code{:165}). The status read lines 1372 to 1379 and not the function they call.

\fhead{Proposed change} Delete the finding (F9) from \code{st:316-318} and the line
\code{st:164}; nothing is to be reported upstream. Correct the memory of the earlier reviews that
cite it.
\end{finding}

% ---------------------------------------------------------------- G6
\begin{finding}{minor}{f:bus-zerocheck-charge}{The zerocheck escape is charged in four different ways, none tight, one to the wrong challenge}
\fhead{Severity} Minor: every version is an upper bound, so none makes a false theorem; the
definition of \lean{busError} is nevertheless undetermined.

\fhead{Classification} An error of the blueprint.

\fhead{Evidence} \code{bp:333} (row \emph{Seams}): ``charged in the bus phase, coordinate by
coordinate as $ζ$ is drawn, $1/|\E|$ per coordinate per constraint''. \code{bp:974}:
``\lean{busError} is $4·2^{μ_{\mathrm{bus}}}/|\E|$ on the $(α, β)$ message plus
\code{gkrError 3 μ_bus}'', no zerocheck term. \code{bp:1001-1003}: ``charged to the $(α, β)$ and GKR
challenges of Layer 6, where $ζ$ is drawn, as an extra $τ_{\max}/|\E|$ per constraint''. The
tracker, in its hole for the bus phase's security (P2): the second plus the first. $(α, β)$ is
drawn before the roots are sent; no coordinate of $ζ$ exists then (\cref{sec:gen-bus-gkr}).

\fhead{Proposed change} One definition, in Layer 6, replacing the three passages:
``\lean{busError}: $4·2^{μ_{\mathrm{bus}}}/|\E|$ on $(α, β)$; on the GKR's challenges,
\lean{gkrError}. The zerocheck escape adds nothing: on each coordinate of $ζ$ (the last layer's
round challenges and its two combination challenges) the state is the conjunction of the GKR's
and of `every constraint's extension, partially evaluated, is zero', and a conjunction escapes
with the larger of the two probabilities, the second being $1/|\E|$ whatever the number of
constraints.'' Row \emph{Seams}: drop ``per constraint''; Layer 7: drop ``the $(α, β)$ and''.
\end{finding}

% ---------------------------------------------------------------- G7
\begin{finding}{minor}{f:bus-leaf-order}{The order of the leaf stacks is nowhere stated, and the spine's \texttt{tuples} has the opposite order}
\fhead{Classification} An error of the blueprint (an omission of a transcribed convention).

\fhead{Evidence} The deployed order puts the three blocks no table owns first, then the tables'
(\cref{sec:gen-bus-leaves}; \code{layout.rs:354-410}; \code{py:507}, ``the blocks no table owns,
stacked first'', \code{py:499}). Ties between equal sizes go by that index, so whenever a table's
height equals $κ_{\mathrm{mem}}$ or $κ_{\mathrm{bc}}$ the order decides the offsets, the
selectors, and so the verdict. The blueprint states the rule for the witness stack only
(\code{bp:322}); for the leaves it cites files (\code{bp:196}). The spine lists the table tuples
first:
\begin{leancode}
-- LeanerVM/Protocol/Spine/Instance.lean:185-186
def tuples (q : Column I.μ) (s : Side) : List (Vector K 16) :=
  I.flushTuples q s ++ I.boundaryTuples q s
\end{leancode}
\src{LeanerVM/Protocol/Spine/Instance.lean:185-186 at b435631}

which is immaterial to \lean{Balanced} (a permutation) and a trap for whoever derives the leaf
layout from it. This is the same kind of finding as the status file's finding on the tie order of
the witness stack (F17).

\fhead{Proposed change} A new row of the pinned conventions: ``Leaf stacks | Per side, the blocks
in index order are the boundary blocks of that side in the order of \code{I.boundary} (state,
memory, bytecode for leanVM), then for each table in order its flushes of that side in order;
the count side has, for each table in order, one block per count column in the order of
\code{I.counts}. Stacked largest first at aligned offsets, ties by that index, no floor on the
depth (\code{layout.rs:354-415}, \code{leaf.rs:149-156}). The depth of the three trees is the push
side's.'' And a guard in Layer 6's tests on sizes with $τ_j = κ_{\mathrm{mem}}$.
\end{finding}

% ---------------------------------------------------------------- G8
\begin{finding}{minor}{f:bus-open-orders}{Orders the blueprint leaves open: the two roots, the boundary evaluations, the linear claims}
\fhead{Classification} An error of the blueprint (omissions of transcribed conventions).

\fhead{Evidence and proposed change} Each a clause to add to Layer 6:
\begin{itemize}
\item the roots message is $(R, R_c)$, the bus root first (\code{gkr.rs:274-275}, \code{py:434}),
  against the specification's prose order;
\item the boundary evaluations are sent in the order in which the sides (push, pull), their blocks
  and their coordinates first name a committed column, a column already valued at the same point
  being skipped (\code{leaf.rs:428-439, 466-471}); for leanVM \code{mem_0}, \code{mem_1},
  \code{mem_2}, \code{cntfin_mem}, \code{cntfin_bc};
\item \code{BusOut.linear} lists the zerocheck claims table by table, then push, pull, count, which
  is the order the table phase gives its powers of $ξ$ (\code{bp:323}; dossier \file{gt-bus.md},
  section E.1);
\item $ζ = (u_0, u_1, χ_0, …, χ_{μ−3})$ of the last layer (\cref{sec:gen-bus-gkr}).
\end{itemize}
\end{finding}

% ---------------------------------------------------------------- G9
\begin{finding}{minor}{f:bus-canonical-checks}{The canonical-encoding checks and the consumption check are absent from the blueprint}
\fhead{Classification} An error of the blueprint (omission).

\fhead{Evidence} The checks of the upper limbs of the announced sizes (B2), of the third limb of
the root's halves (B10) and of the consumption of the stream (B17), in \cref{tab:gen-bus-checks}.
No passage of the blueprint concerns canonical encodings or the consumption of the stream: a
search for ``canonical'' finds only \code{bp:1294}, the public word of acceptance test 10;
``upper limb'' finds nothing; ``consumed'' occurs only in the sense of ``used by'' (for example
\code{bp:129, 146, 262}). (The citation check corrected the wording: the dossier said that the
search for the three words finds only the public word of acceptance test 10.) The status lists
the Rust's rejections as ``four predicates plus truncations'' (\code{st:319-320}), without
\code{NonCanonicalEncoding} and \code{NotFullyConsumed}. None protects soundness; without them
\lean{verify} accepts streams the Rust rejects, and \lean{verify} is specified as the verifier
``the Rust prover's proofs are checked against''. Layer 12's six mutations (\code{bp:1238-1239})
exercise none.

\fhead{Proposed change} Layer 12: list the three checks among \lean{verify}'s, with their source
lines; add three mutations (a nonzero upper limb in an announced size, a nonzero third limb in a
half of the root, one element appended to the stream).
\end{finding}

% ---------------------------------------------------------------- G10
\begin{finding}{minor}{f:bus-side-conditions}{Two side conditions a bus phase over an abstract instance needs}
\fhead{Classification} An error of the blueprint (the spine's instance admits instances no bus
phase of the deployed shape serves; leanVM's is not one).

\fhead{Evidence} (The dossier's probe on the spine, section H.2, parts 1 and 2. The citation
check could not re-run it, Lean not running in the checkout since the upgrade; the probe's source
is identical to the dossier's reproduction.)
\begin{enumerate}
\item With $d = 0$ and a count column, \lean{M3Holds} is inhabited and the count form's one
  polynomial, $X_c$, has total degree 1: a statement carrying it is outside \lean{Seam.bus}, whose
  second conjunct is \code{t.poly.totalDegree ≤ I.d} (\code{Seams.lean:176}). Completeness of a
  bus phase that emits the count form fails on that instance.
\item A table with a constraint and no flush, no count column and no boundary block of its height
  is not on the bus: the trees have depth 0, $ζ$ has no coordinate, and the table's zerocheck
  point needs $τ_j$. The Rust guards the inequality at the start of the table sumcheck
  (\code{constraints.rs:250-253}).
\end{enumerate}

\fhead{Proposed change} Either two fields of \lean{M3Instance}, \code{one_le_d : 1 ≤ d} and
\code{constrained_on_bus : ∀ j, constraints j ≠ [] → τ j ≤ μ_bus}, or two hypotheses of the bus
phase, stated in Layer 6; and acceptance test 28 to name the first.
\end{finding}

% ---------------------------------------------------------------- G11
\begin{finding}{minor}{f:bus-stale-sketch}{The sketch of Layer 6 and the tracker's bus section are stale, and the invariant they name is not one}
\fhead{Classification} An error of the blueprint (text the spine superseded).

\fhead{Evidence} \code{bp:954-979} has \code{StmtIn := Unit}, a \code{BusOut} with fields
\code{ζ, rem, pool, α, β}, a \code{relOut} with ``count root $≠ 0$'' and without the zerocheck
claims, the public lines and \code{aux}, and a \code{LeafLayout prog s} that nothing defines; the
spine has \code{I.Stmt}, \code{BusOut} with \code{linear} and \code{columns}, and \lean{Seam.bus}.
The tracker's hole for the bus phase (P1) has
\code{Phase.Def I Unit (BusOut I) (busSpec I) (M3Holds I) (Seam.bus I)} and cites
``specification (5.4)'', which does not exist (\code{st:303}). \code{bp:975-977}: ``The knowledge
state function's invariant after $(α, β)$ is `the pushed and pulled multisets of $q$ differ, or
some count is zero, or a boundary claim is false'{}''. After $(α, β)$ the first clause must be
``the two products at $(α, β)$ differ'': if the multisets differ and the products collide, the
prover is honest from there on and no later challenge can be charged. And no boundary claim
exists before the last message.

\fhead{Proposed change} Replace the code block of Layer 6 by the slot's types (dossier
\file{gt-bus.md}, section C) and the schedule
\code{V_to_P (α, β) ; P_to_V (R, R_c) ; the GKR's ; P_to_V the boundary evaluations}; replace the
sentence on the invariant by the table of the dossier's section D.3.
\end{finding}

% ---------------------------------------------------------------- G12
\begin{finding}{minor}{f:bus-binary-combination}{The binary layer's combination challenge has no error assigned}
\fhead{Classification} An error of the blueprint (omission).

\fhead{Evidence} \lean{gkrError} (\code{bp:941-942}) names combiners, rounds and pairs; for odd
$μ$ the first layer has one combination challenge (\code{gkr.rs:387}), whose error is $1/|\E|$.

\fhead{Proposed change} ``$1/|\E|$ to each combination challenge (two per radix-4 layer, one for
the binary layer)''.
\end{finding}

% ---------------------------------------------------------------- G13
\begin{finding}{minor}{f:bus-status-one-root}{The status file's finding on the one bus root misreads the pinned specification}
\fhead{Classification} An error of the status.

\fhead{Evidence} The status file's finding on the specification's listing of the bus roots (S11),
\code{st:289-290}: ``§8.5 lists the bus roots as `the count root $R_c$ and one bus root $R$' but
does not say the push and pull roots are one scalar''. It does: ``one $R$ for both sides is the
balance check'' (\code{08:69}) and ``the one bus root for both sides'' (\code{08:100}). The status
file's finding ``one root'' (F3) is an agreement of the three sources, not a divergence.

\fhead{Proposed change} Delete the finding on the listing of the bus roots (S11); reword the
finding ``one root'' (F3) as a convention.
\end{finding}

% ---------------------------------------------------------------- G14
\begin{finding}{minor}{f:bus-citation-drift}{Citation drift}
\fhead{Evidence} \code{bp:197} cites \code{leaf.rs:664-668} for the count blocks: those lines are
the signature of \code{prove_balance}; the count side is \code{leaf.rs:606-622}. \code{bp:601} has
``Lemma 5.1'' for the multiset lemma, which is Lemma 5.2 (Theorem 5.1 is the product check,
\code{05:36-48}), as Layer 5 itself writes. The tracker's hole for the bus phase (P1) has
``specification (5.4)''.
\end{finding}

% ---------------------------------------------------------------- G15
\begin{finding}{note}{f:bus-disagreements}{Disagreements between the three sources (seven, none a divergence of the two verifiers)}
The blueprint follows the Rust on each, rightly. The seven are items 1 to 7 of
\cref{tab:gen-bus-disagreements}: the announced rate; the order of the roots; the seed (the
status file's finding on the seed, F14); the checks of the verifiers the specification does not
mention; the combiner after the last layer (\cref{f:bus-extra-combiner}); the program as a stacked
table in the Python verifier; the fingerprint bound, the count columns and the Rust's layout
assertions. (Item 8 of the table was added by the citation check.)
\end{finding}

% ---------------------------------------------------------------- G16
\begin{finding}{note}{f:bus-table-phase-tables}{The table phase must know which tables it opens}
For the dossier on the table sumcheck. The deployed sumcheck has $τ_{\max} = \max$ over the six
tables (\code{constraints.rs:250}, \code{py:607}) and opens every column of those six and no other
(\code{mod.rs:344-351}). In the instance the shared columns are tables too, of log-height
$κ_{\mathrm{mem}} ≥ 16$; a table phase taking its number of rounds from every table of \code{I}
has a different transcript. The criterion ``has a constraint, a flush or a count column'' is
derivable and must be written down.
\end{finding}

% ---------------------------------------------------------------- G17
\begin{finding}{note}{f:bus-no-rbr-in-spec}{The specification gives no round-by-round analysis of the bus phase}
(Dossier \file{gt-bus.md}, section D.1.) \lean{gkrError} and the zerocheck charge are the
blueprint's own analysis; the specification cannot be cited for them, and the Rust asserts only a
coarse sum.
\end{finding}
